% !TEX program = xelatex
% !TEX options = -synctex=1 -interaction=nonstopmode -file-line-error
% example-maths.tex
% Demonstrates rossbeamermaths with chalkboard option: theorem, lemma,
% proof, definition, \eqhl, \step.
\documentclass[palette=blackboard, chalkboard]{rossbeamermaths}

\title{Introduction to Measure Theory}
\subtitle{Lebesgue Integration on $\mathbb{R}^n$}
\author{Dr.~Eleanor Hartley}
\institute{Department of Pure Mathematics, University of Oxford}
\date{Hilary Term 2026}

\begin{document}

% --- Title ---
\begin{frame}[plain]
  \titlepage
\end{frame}

% --- Definition ---
\begin{frame}{Sigma Algebras}
  \begin{definition}[$\sigma$-algebra]
    Let $X$ be a non-empty set. A collection $\mathcal{F} \subseteq 2^X$
    is a \textbf{$\sigma$-algebra} if:
    \begin{enumerate}
      \item $X \in \mathcal{F}$
      \item $A \in \mathcal{F} \implies A^c \in \mathcal{F}$
      \item $A_1, A_2, \ldots \in \mathcal{F}
            \implies \bigcup_{n=1}^{\infty} A_n \in \mathcal{F}$
    \end{enumerate}
  \end{definition}
  \vskip6pt
  The pair $(X, \mathcal{F})$ is called a \textbf{measurable space}.
\end{frame}

% --- Theorem ---
\begin{frame}{Monotone Convergence}
  \begin{theorem}[Monotone Convergence Theorem]
    Let $(f_n)$ be a sequence of measurable functions with
    $0 \le f_1 \le f_2 \le \cdots$ pointwise. Then
    \[
      \int \lim_{n \to \infty} f_n \, d\mu
      = \lim_{n \to \infty} \int f_n \, d\mu.
    \]
  \end{theorem}
\end{frame}

% --- Lemma with eqhl ---
\begin{frame}{Supporting Lemma}
  \begin{lemma}[Fatou]
    If $(f_n)$ is a sequence of non-negative measurable functions, then
    \[
      \int \liminf_{n \to \infty} f_n \, d\mu
      \;\le\;
      \eqhl{\liminf_{n \to \infty} \int f_n \, d\mu}.
    \]
  \end{lemma}
  \vskip6pt
  The highlighted term shows that the limit inferior and integral
  do not freely commute without additional hypotheses.
\end{frame}

% --- Proof with steps ---
\begin{frame}{Proof Sketch}
  \begin{proof}[Proof of the Monotone Convergence Theorem]
    \resetsteps
    \step{Define $f = \lim_{n \to \infty} f_n$, which exists pointwise
          since $(f_n)$ is increasing.}
    \step{Since $f_n \le f$ for all $n$, we have
          $\int f_n \, d\mu \le \int f \, d\mu$, giving
          $\limsup \int f_n \le \int f$.}
    \step{For the reverse inequality, fix a simple function
          $0 \le s \le f$ and $\alpha \in (0,1)$. Define
          $E_n = \{x : f_n(x) \ge \alpha \, s(x)\}$.}
    \step{Then $E_n \uparrow X$ and
          $\int f_n \ge \alpha \int_{E_n} s \, d\mu
           \to \alpha \int s \, d\mu$.}
    \step{Taking $\sup$ over all simple $s \le f$ and letting
          $\alpha \to 1$ gives
          $\liminf \int f_n \ge \int f$.}
  \end{proof}
\end{frame}

% --- Equation box ---
\begin{frame}{Key Inequalities}
  \begin{eqbox}[Dominated Convergence]
    If $|f_n| \le g$ with $g \in L^1(\mu)$ and $f_n \to f$ a.e., then
    \[
      \eqhl[rbChalkBlue]{\lim_{n \to \infty} \int |f_n - f| \, d\mu = 0}
    \]
  \end{eqbox}
  \vskip6pt
  This is the workhorse theorem for exchanging limits and integrals
  in applications.
\end{frame}

% --- Conclusion ---
\conclusionslide{Questions?}{e.hartley@maths.ox.ac.uk}

\end{document}
